\chapter{三角形境界要素による静電場}
\label{ch:electrostatics}

\section{P0 表面電荷の積分}
自由空間の Coulomb 定数を $\kc=(4\pi\eps)^{-1}$ とすると、
式\eqref{eq:p0}の source が作る電位と電場は
\begin{align}
 \phi(\vect x)&=\kc\sum_{i=1}^{N}\frac{q_i}{A_i}
       \int_{T_i}\frac{1}{\lVert\vect x-\vect y\rVert}\dd S_{\vect y},
       \label{eq:panel_phi}\\
 \vect E_{\mathrm s}(\vect x)&=\kc\sum_{i=1}^{N}\frac{q_i}{A_i}
       \int_{T_i}\frac{\vect x-\vect y}
       {\lVert\vect x-\vect y\rVert^3}\dd S_{\vect y}.
       \label{eq:panel_field}
\end{align}
BEACH の公開 source kernel は \code{triangle_p0} に固定されている。
要素電荷は既知の蓄積量として与えられるため、絶縁体の通常経路では、
境界電位から全 $q_i$ を毎回逆算する連立方程式は解かない。
浮遊導体の等電位化だけが、追加の境界条件を連立方程式へ課す経路である。

P0 は要素ごとの一定密度近似であり、真の表面電荷密度の連続性や細部を完全には表さない。
面の積分を正確に評価しても、三角形形状と密度の離散化誤差は残る。
面上の一定・線形 source に関する解析積分の系譜は Graglia の研究に見られる\citep{Graglia1993}。

\section{辺の対数項と立体角}
解析 kernel の構造を一つの三角形について示す。評価点 $\vect x$ の面からの符号付き高さを
$d$、平面への投影を $\vect x_\parallel$、辺 $e$ の面内外向き法線を $\vect\nu_e$ とする。
辺の長さを $\ell_e$、両端までの距離を $r_{e,1},r_{e,2}$ とおくと、
\begin{equation}
 L_e=\log\frac{r_{e,1}+r_{e,2}+\ell_e}
                   {r_{e,1}+r_{e,2}-\ell_e},\qquad
 d_e=(\vect a_e-\vect x_\parallel)\cdot\vect\nu_e
\end{equation}
を用いて
\begin{align}
 I_\phi&=\sum_e d_eL_e-d\,\Omega,\\
 \vect I_E&=\sum_e\vect\nu_eL_e+\vect n\Omega,\\
 \phi_i&=\kc\frac{q_i}{A_i}I_\phi,\qquad
 \vect E_i=\kc\frac{q_i}{A_i}\vect I_E
\end{align}
を評価する。$\Omega$ は実装の法線規約に従う符号付き立体角であり、
\src{bem_panel_kernel.f90} では三頂点の相対ベクトルの三重積と内積から
\code{atan2} で計算する。
上式は面外での構造を表し、面上では専用の自己項と trace を使う。

重心に点電荷を置く近似と異なり、三角形の広がりを近傍でも保持する。
公開入力に経験的な softening 長を導入して特異性を置き換えるモデルではない。
鋭い辺や頂点の近傍では、電場の特異な振る舞いと数値評価規約を区別する必要がある。

\section{自己電位と片側電場}
要素内の通常の評価点では、面上の電位は連続であり、面積を持つ三角形の自己電位も有限である。
法線電場は表面電荷によって跳ぶ。
\begin{equation}
 \vect E_i^{\pm}
 =\vect E_i^{\mathrm{PV}}\pm\frac{\sigma_i}{2\eps}\vect n_i,
 \qquad
 \vect n_i\cdot(\vect E_i^+-\vect E_i^-)=\frac{\sigma_i}{\eps}.
 \label{eq:jump}
\end{equation}
$\mathrm{PV}$ は両側極限の平均としての principal-value trace である。
どちらを真空・プラズマ側として採用するかは \code{surface_side} と法線の向きで決める。
\code{normal_plus}、\code{normal_minus}、整合した閉曲面の \code{outward_closed} は、
幾何の向きとセットで解釈する。

衝突判定は両側から行われるため、衝突の表裏と電場の片側 trace は同じ概念ではない。
また、面上ちょうどの Direct と FMM には trace 規約の差があるため、
通常の solver 一致試験は面から離した点、または規約を揃えた自己項で行う。

\section{内部の長さ正規化と SI 単位}
代表長 $L$ に対して $\vect x=L\widehat{\vect x}$ と変換した場合、
要素総電荷 $q_i$ を保持したまま、電位は $L^{-1}$、電場は $L^{-2}$ の係数で SI へ戻る。
\begin{equation}
 \phi=\frac{\kc}{L}\widehat\phi,\qquad
 \vect E=\frac{\kc}{L^2}\widehat{\vect E}.
\end{equation}
正規化は数値尺度を整える処理であり、物理的な遮蔽や mesh 細分化を代替しない。
高次 moment は長さの次数が異なるため、周期 operator の生成でも尺度の揃え方が精度に影響する。

\basisdoc{\src{docs/DirectSolver.md}、\src{src/mesh/panel/bem_panel_geometry.f90}、
\src{src/physics/field_solver/panel/bem_panel_kernel.f90}、
\src{src/physics/field_solver/panel/bem_panel_self_terms.f90}。}
